\documentclass[11pt]{article}

\usepackage[margin=1in]{geometry}
\usepackage{amsmath,amssymb,amsthm,mathtools}
\usepackage{booktabs,longtable,array}
\usepackage{enumitem}
\usepackage{xcolor}
\usepackage{hyperref}
\usepackage{listings}
\usepackage[T1]{fontenc}
\usepackage[utf8]{inputenc}

\hypersetup{
  colorlinks=true,
  linkcolor=blue!50!black,
  citecolor=blue!50!black,
  urlcolor=blue!50!black
}

\lstdefinestyle{cstyle}{
  language=C,
  basicstyle=\ttfamily\small,
  keywordstyle=\color{blue!60!black},
  commentstyle=\color{green!40!black},
  stringstyle=\color{orange!60!black},
  breaklines=true,
  frame=single,
  columns=fullflexible
}
\lstset{style=cstyle}

\newcommand{\Z}{\mathbb Z}
\newcommand{\Q}{\mathbb Q}
\newcommand{\R}{\mathbb R}
\newcommand{\E}{\mathbb E}
\newcommand{\Prob}{\mathbb P}
\newcommand{\one}{\mathbf 1}
\newcommand{\KL}{D_{\mathrm{KL}}}
\newcommand{\Htwo}{H_2}
\newcommand{\Ccrit}{C_6}
\newcommand{\ord}{v_2}

\title{Specifications for C Computation Tests:\\
Entropy-Space, Doob-Transform, Schr\"odinger-Bridge, and Transfer-Operator Diagnostics for the Accelerated Collatz Cocycle}
\author{Research notes and implementation specification}
\date{\today}

\begin{document}
\maketitle

\begin{abstract}
This document specifies a C-based computational program for testing four related formulations of the remaining Collatz obstruction: the Doob $h$-transform/Cram\'er tilt, Schr\"odinger bridge/entropic interpolation, deterministic transport between annealed height laws, and transfer-operator flow over residue-height states. The computations are not intended to prove Collatz. They are designed to test whether the deterministic accelerated Collatz orbit shadows, avoids, or is repelled from the annealed rare-event laws predicted by the forward height cocycle. The central empirical object is the valuation cocycle
\[
  b(x)=\ord(3x+1),\qquad \Delta_b=\log_2 3-b,
\]
with annealed law $p_b=2^{-b}$ and moment transform
\[
  M(s)=\sum_{b\ge1}2^{-b}2^{s(\log_2 3-b)}=\frac{3^s}{2^{1+s}-1}.
\]
At the Cram\'er root $s=1$, the tilted law is
\[
  p_b^{(1)}=p_b2^{\Delta_b}=3\cdot4^{-b}=\frac34\left(\frac14\right)^{b-1}.
\]
The tests below are designed to determine how critical excursions, post-exit floors, and ordinary orbit windows sit relative to this tilted law and to the associated transfer-operator spectra.
\end{abstract}

\tableofcontents
\newpage

\section{Executive Summary}

The program has converged on a single analytic obstruction: the annealed height law is explicit and correctly predicts rare-event statistics, while the needed theorem is quenched and pointwise. In computational terms, the question becomes:
\begin{quote}
Can a deterministic positive integer orbit repeatedly shadow the annealed rare-event flow, or does doing so force a finite spectral/character obstruction already ruled out by the finite certificate machinery?
\end{quote}

The present specification builds computational tests around four formulations:
\begin{enumerate}[label=\textbf{F\arabic*.},leftmargin=2.5em]
  \item \textbf{Doob $h$-transform / Cram\'er tilt.} Critical excursions should have empirical valuation law close to $p^{(1)}_b=3\cdot4^{-b}$ in the bulk. Post-exit segments should fail to reacquire this law if the taboo mechanism is real.
  \item \textbf{Schr\"odinger bridge / entropic interpolation.} A critical excursion is a finite bridge: entry boundary layer, tilted bulk, exit boundary layer. Consecutive critical bridges should be forbidden or spectrally suppressed.
  \item \textbf{Deterministic transport in entropy space.} Sliding windows of valuation words trace paths in the simplex of empirical valuation laws. Ordinary windows should lie near $p$, high ascents near $p^{(1)}$, and post-exit windows should be separated from $p^{(1)}$.
  \item \textbf{Transfer-operator flow over residue-height states.} A family $\mathcal L_s$ over residue-height states should encode ordinary dynamics ($s=0$), critical tilt ($s=1$), QSD, H23a character gaps, and post-exit taboo through spectra such as $\rho(\mathcal L_1\mathcal E\mathcal L_1)$.
\end{enumerate}

The outputs are standardized CSV/JSON logs and summary tables. All exact logical predicates, especially criticality inequalities, must be evaluated using integer arithmetic. Floating-point arithmetic is used only for diagnostics such as entropy, pressure, and KL divergence.

\section{Mathematical Definitions}

\subsection{Accelerated odd map and valuation word}

For odd positive $x$, define
\[
  S(x)=\frac{3x+1}{2^{b(x)}},\qquad b(x)=\ord(3x+1).
\]
An orbit segment of length $m$ has valuation word
\[
  w=(b_0,b_1,\ldots,b_{m-1}),
  \qquad b_i=b(S^i x),
\]
and cumulative valuation
\[
  D_m=\sum_{i=0}^{m-1}b_i.
\]
The deficit/excess relative to the equilibrium slope is
\[
  E_m=D_m-m\log_2 3.
\]
In the forward height convention, define
\[
  \Delta_b=\log_2 3-b,
  \qquad Y_m=\sum_{i=0}^{m-1}\Delta_{b_i}=m\log_2 3-D_m=-E_m.
\]
High ascents correspond to large positive $Y_m$.

\subsection{Annealed valuation law and pressure transform}

Under the Haar/annealed model,
\[
  p_b=\Prob(b=k)=2^{-b},\qquad b\ge1.
\]
The base-$2$ moment transform of the height increment is
\[
  M(s)=\E_p\bigl(2^{s\Delta_b}\bigr)
  =\sum_{b\ge1}2^{-b}2^{s(\log_2 3-b)}
  =\frac{3^s}{2^{1+s}-1}.
\]
The Cram\'er root is $s^*=1$, since $M(1)=1$. The tilted law at parameter $s$ is
\[
  p_b^{(s)}=\frac{p_b2^{s\Delta_b}}{M(s)}.
\]
At $s=1$,
\[
  p_b^{(1)}=3\cdot4^{-b}=\frac34\left(\frac14\right)^{b-1}.
\]
The binary entropy associated with the critical bulk law is
\[
  \Htwo(3/4)=-\frac34\log_2\frac34-\frac14\log_2\frac14\approx0.811278.
\]

\subsection{Critical excursions}

Let $V$ be an odd launch floor and $P$ the peak reached during the forward excursion before the first return below $V$ or before a specified stopping condition. Define the logarithmic height
\[
  h(V)=\log_2\frac{P}{V}.
\]
The $A$-critical launch set is
\[
  C_A=\{V:\ h(V)\ge \log_2 V-A\}.
\]
Equivalently, for $A=6$,
\[
  V\in C_6\quad\Longleftrightarrow\quad V^2\le 2^6P.
\]
This predicate must be evaluated exactly as an integer inequality.

\subsection{Empirical valuation laws}

For a finite word $w$ of length $m$, define its empirical valuation law
\[
  \mu_w(b)=\frac{1}{m}\#\{0\le i<m:b_i=b\}.
\]
Diagnostics include
\[
  \bar b(\mu)=\sum_b b\mu(b),\qquad
  \bar\Delta(\mu)=\sum_b(\log_2 3-b)\mu(b),
\]
\[
  H(\mu)=-\sum_b\mu(b)\log_2\mu(b),
\]
\[
  \KL(\mu\|p)=\sum_b\mu(b)\log\frac{\mu(b)}{p_b},
  \qquad
  \KL(\mu\|p^{(1)})=\sum_b\mu(b)\log\frac{\mu(b)}{p_b^{(1)}}.
\]

\section{C Implementation Architecture}

\subsection{Directory layout}

Recommended structure:
\begin{verbatim}
collatz_entropy_tests/
  include/
    collatz_u128.h
    valuation_word.h
    histograms.h
    entropy_stats.h
    csv_writer.h
    operator_matrix.h
    config.h
  src/
    collatz_u128.c
    valuation_word.c
    histograms.c
    entropy_stats.c
    csv_writer.c
    operator_matrix.c
    common_cli.c
  tools/
    doob_tilt_test.c
    bridge_layers_test.c
    entropy_transport_test.c
    transfer_operator_test.c
    critical_floors_test.c
    post_exit_test.c
    pressure_tail_test.c
  scripts/
    build.sh
    run_all.sh
    summarize.py
  data/
    raw/
    summaries/
  docs/
\end{verbatim}

\subsection{Compiler and build flags}

Use a strict, reproducible build:
\begin{lstlisting}
cc -O3 -std=c17 -Wall -Wextra -Wshadow -Wconversion \
   -pedantic -march=native -DNDEBUG \
   -Iinclude src/*.c tools/doob_tilt_test.c -lm -o bin/doob_tilt_test
\end{lstlisting}

For debug builds:
\begin{lstlisting}
cc -O0 -g -fsanitize=address,undefined -std=c17 \
   -Wall -Wextra -Wshadow -Wconversion -pedantic \
   -Iinclude src/*.c tools/doob_tilt_test.c -lm -o bin/doob_tilt_test_dbg
\end{lstlisting}

\subsection{Integer arithmetic requirements}

Use unsigned 128-bit arithmetic for all exact checks up to the intended bound:
\[
  \texttt{\_\_uint128\_t}.
\]
Criticality must be checked as
\[
  V^2\le 2^A P
\]
without floating point. The code must detect overflow before multiplication. If computations exceed $128$ bits, either stop with a clear diagnostic or switch to a big-integer backend such as GMP. The first implementation target should remain within the safe $128$-bit range.

\subsection{Common command-line options}

Every tool should support:
\begin{verbatim}
--start <uint64>          first odd V or x to test
--end <uint64>            exclusive upper bound
--octave-min <N>          first dyadic octave
--octave-max <N>          last dyadic octave
--A <int>                 criticality parameter, default 6
--max-steps <int>         maximum accelerated steps per orbit/excursion
--max-b <int>             histogram cap for valuation b
--window <int>            sliding window length
--out <path>              output CSV path
--summary <path>          output summary CSV/JSON path
--seed <uint64>           optional deterministic sample seed
--sample-rate <int>       test every kth odd integer, default 1
--threads <int>           optional threading
--verbose                 progress and diagnostics
\end{verbatim}

\subsection{Shared data structures}

\begin{lstlisting}
typedef struct {
    uint64_t start;
    uint64_t end;
    int A;
    int max_steps;
    int max_b;
    int window;
    const char *out_path;
    const char *summary_path;
} RunConfig;

typedef struct {
    uint64_t V;          // launch floor
    uint64_t exit_floor; // first return below V, if any
    __uint128_t peak;    // peak P
    int steps;
    int max_b;
    int *b_word;         // length steps
    double height;       // diagnostic only
    double amin;         // log2(V^2/P), diagnostic only
    bool critical;
} ExcursionRecord;

typedef struct {
    uint64_t count;
    uint64_t *bins;      // bins[1..max_b], overflow bin max_b+1
    uint64_t overflow;
} ValuationHist;
\end{lstlisting}

\section{Core Algorithms}

\subsection{Accelerated step}

For odd $x$, compute
\[
  y=3x+1,\\qquad b=\ord(y),\qquad S(x)=y/2^b.
\]
The implementation must check overflow in $3x+1$.

\begin{lstlisting}
static inline int ctz_u128(__uint128_t y) {
    int c = 0;
    while ((y & 1u) == 0u) { y >>= 1; c++; }
    return c;
}

bool accelerated_step_u128(__uint128_t x, __uint128_t *next, int *b) {
    __uint128_t y = 3*x + 1;
    *b = ctz_u128(y);
    *next = y >> (*b);
    return true;
}
\end{lstlisting}

\subsection{Excursion extraction}

Given an odd launch floor $V$, iterate until the first return below $V$, a maximum step bound, overflow, or terminal cycle detection.
Record:
\begin{itemize}
  \item valuation word $b_0,\dots,b_{m-1}$;
  \item peak $P=\max_i x_i$;
  \item height $h=\log_2(P/V)$;
  \item criticality $V^2\le2^A P$;
  \item exit floor $E(V)$ if the excursion returns below $V$.
\end{itemize}

\subsection{Exact criticality check}

For $A=6$, test
\[
  V^2\le 64P.
\]
General $A$:
\[
  V^2\le 2^A P.
\]

\begin{lstlisting}
bool is_A_critical_u128(uint64_t V, __uint128_t P, int A) {
    __uint128_t lhs = (__uint128_t)V * (__uint128_t)V;
    __uint128_t rhs = P << A;
    return lhs <= rhs;
}
\end{lstlisting}

\subsection{Histogram and KL divergence}

For each word $w$, build a histogram over $b\in\{1,\dots,B_{\max}\}$, with overflow bin. Compute $\mu_w$. Compare to
\[
  p_b=2^{-b},\qquad p_b^{(1)}=3\cdot4^{-b}.
\]
For the overflow bin, use the exact tail mass:
\[
  p_{>B}=2^{-B},
\]
for the annealed law, and
\[
  p^{(1)}_{>B}=4^{-B},
\]
for the tilted law, since $p_b^{(1)}=3\cdot4^{-b}$.

\section{Test Group I: Doob $h$-Transform and Cram\'er Tilt}

\subsection{Test I.1: Critical-excursion tilted valuation law}

\paragraph{Goal.} Confirm that the bulk valuation law of $A$-critical excursions converges to
\[
  p_b^{(1)}=3\cdot4^{-b}.
\]

\paragraph{Input.} Odd launch floors $V\in[2^N,2^{N+1})$ for octaves $N_{\min}\le N\le N_{\max}$.

\paragraph{Procedure.}
\begin{enumerate}
  \item Extract the first-return excursion from each $V$.
  \item Retain only records satisfying $V^2\le2^A P$.
  \item Optionally remove an entry boundary layer of length $r$ and exit boundary layer of length $r$.
  \item Build aggregate valuation histograms for full word and bulk word.
  \item Compute $\widehat p_b$, $\KL(\widehat p\|p)$, $\KL(\widehat p\|p^{(1)})$, entropy, and mean height increment.
\end{enumerate}

\paragraph{Output CSV.}
\begin{verbatim}
octave,A,count_critical,total_words,total_symbols,
b,p_hat,p_annealed,p_tilted,kl_to_p,kl_to_tilt,entropy,mean_b,mean_delta
\end{verbatim}

\paragraph{Expected signal.}
For critical-excursion bulk symbols,
\[
  \widehat p_b\to3\cdot4^{-b},\qquad
  \KL(\widehat p\|p^{(1)})\to0.
\]

\subsection{Test I.2: Height-tail law by octave}

\paragraph{Goal.} Verify octave-uniform tail
\[
  \Prob(h\ge H)\asymp2^{-H},
\]
and therefore
\[
  \#(C_A\cap[2^N,2^{N+1}))=O_A(1).
\]

\paragraph{Procedure.}
For each octave $N$ and threshold $H$, count
\[
  \#\{V\in[2^N,2^{N+1}):h(V)\ge H\}.
\]
Fit
\[
  \log_2 \Prob_N(h\ge H)\approx -H+c.
\]

\paragraph{Output CSV.}
\begin{verbatim}
octave,H,count,total,prob,log2_prob,slope_fit,intercept_fit
\end{verbatim}

\paragraph{Expected signal.}
For large $H$ below finite-size cutoff,
\[
  \log_2\Prob_N(h\ge H)+H
\]
should be approximately constant across octaves.

\subsection{Test I.3: Post-exit conditional valuation law}

\paragraph{Goal.} Determine whether post-exit segments are anti-tilted away from $p^{(1)}$.

\paragraph{Procedure.}
For each $V\in C_A$ with exit floor $E(V)$:
\begin{enumerate}
  \item Start a new segment at $E(V)$.
  \item Record the first $L_{post}$ valuations or the next first-return excursion.
  \item Build the empirical law $\mu_{post}$.
  \item Compute $\KL(\mu_{post}\|p^{(1)})$ and empirical post pressure
  \[
    M_{post}(1)=\sum_b\mu_{post}(b)2^{\log_2 3-b}.
  \]
\end{enumerate}

\paragraph{Expected signal.}
If post-exit taboo is pressure-based, then
\[
  M_{post}(1)<1
\]
for the relevant initial post-exit window, or at least
\[
  \KL(\mu_{post}\|p^{(1)})\ge c>0.
\]

\section{Test Group II: Schr\"odinger Bridge / Entropic Interpolation}

\subsection{Test II.1: Entry, bulk, and exit boundary layers}

\paragraph{Goal.} Determine whether critical excursions decompose into entry layer, tilted bulk, and exit layer.

\paragraph{Procedure.}
For each critical word $w$ of length $m$, choose boundary depth $r$ and split:
\[
  w=w_{entry}\,w_{bulk}\,w_{exit}.
\]
Compute valuation laws
\[
  \mu_{entry},\quad \mu_{bulk},\quad \mu_{exit}.
\]
Compare each to $p$ and $p^{(1)}$.

\paragraph{Output CSV.}
\begin{verbatim}
octave,V,word_len,r,region,count_symbols,mean_b,mean_delta,
entropy,kl_to_p,kl_to_tilt,p1_freq,p2_freq,p3_freq,overflow
\end{verbatim}

\paragraph{Expected signal.}
The bulk should be closest to $p^{(1)}$. Entry and exit layers should show systematic deviations. The earlier transient dimension effects should be visible in the entry layer.

\subsection{Test II.2: Critical bridge non-concatenation}

\paragraph{Goal.} Test whether the exit marginal of a critical bridge is separated from the entrance marginal of another critical bridge.

\paragraph{Procedure.}
For each critical excursion:
\begin{enumerate}
  \item Record exit floor $E(V)$.
  \item Compute whether $E(V)\in C_A$ exactly.
  \item Compute next-excursion height $h(E(V))$.
  \item Define margin
  \[
    \Delta_{post}=\log_2E(V)-A-h(E(V)).
  \]
\end{enumerate}

\paragraph{Output CSV.}
\begin{verbatim}
V,peak,exit_floor,A,h_V,h_exit,amin_V,delta_post,
critical_again,next_word_len,next_peak
\end{verbatim}

\paragraph{Expected signal.}
For $A=6$, computations so far suggest
\[
  E(C_6)\cap C_6=\varnothing.
\]
A positive lower bound for $\Delta_{post}$ would be stronger.

\subsection{Test II.3: Empirical bridge action}

\paragraph{Goal.} Verify that critical excursions minimize an entropy action predicted by the annealed bridge.

\paragraph{Definition.}
For word $w$ of length $m$, define
\[
  \mathcal A(w)=m\KL(\mu_w\|p).
\]
For a high ascent of height $H$, compare $\mathcal A(w)$ to $H\log 2$.

\paragraph{Expected signal.}
Critical excursions should have action near the Cram\'er prediction. Failed second excursions after exit should show excess action or inability to approach the bridge law.

\section{Test Group III: Deterministic Transport in Entropy Space}

\subsection{Test III.1: Sliding-window empirical-law trajectories}

\paragraph{Goal.} View deterministic orbits as paths in the probability simplex of valuation laws.

\paragraph{Procedure.}
For each orbit and window length $L_w$:
\begin{enumerate}
  \item Compute sliding-window law $\mu_{n,L_w}$.
  \item Record coordinates
  \[
    \bar b,\\quad \bar\Delta,
    \quad H(\mu),
    \quad \KL(\mu\|p),
    \quad \KL(\mu\|p^{(1)}).
  \]
  \item Label windows as ordinary, pre-critical, critical, post-exit, or terminal.
\end{enumerate}

\paragraph{Output CSV.}
\begin{verbatim}
x0,n,window,label,mean_b,mean_delta,entropy,
kl_to_p,kl_to_tilt,height_start,height_end,max_height
\end{verbatim}

\paragraph{Expected signal.}
Ordinary windows cluster near $p$. Critical windows approach $p^{(1)}$. Post-exit windows avoid $p^{(1)}$ if taboo is real.

\subsection{Test III.2: Distance to the Cram\'er curve}

\paragraph{Goal.} Determine whether orbit windows align with the exponential family
\[
  p^{(s)}_b=\frac{p_b2^{s\Delta_b}}{M(s)}.
\]

\paragraph{Procedure.}
For each empirical law $\mu$, find
\[
  s^*(\mu)=\arg\min_{s\in[s_{min},s_{max}]}\KL(\mu\|p^{(s)}).
\]
Record the minimizing $s$ and the residual distance.

\paragraph{Expected signal.}
Critical windows should have $s^*(\mu)\approx1$. Post-exit windows should move away from $s=1$.

\subsection{Test III.3: Entropic action budget}

\paragraph{Goal.} Test whether repeated critical events fail because the post-exit segment lacks enough entropy/height budget to return to the critical entropy manifold.

\paragraph{Procedure.}
For each critical excursion and following segment, compute cumulative action
\[
  \mathcal A_{0:n}=\sum_{j}\KL(\mu_{j}\|p)
\]
or block action
\[
  |w|\KL(\mu_w\|p).
\]
Compare against achieved height.

\paragraph{Expected signal.}
If an action barrier exists, post-exit attempted second ascents should require action greater than the available large-deviation budget.

\section{Test Group IV: Transfer-Operator Flow over Residue-Height States}

\subsection{State space and operators}

For fixed modulus $Q$ and height strip $[0,L]$, define finite bins
\[
  \mathcal X_{Q,L}=\Omega_Q\times\{0,1,\ldots,L\}.
\]
The tilted transfer operator is
\[
  (\mathcal L_s f)(a,z)=\sum_b p_b2^{s(\log_2 3-b)}f(T_ba,z+\log_2 3-b),
\]
with killing at boundaries. In implementation, $z$ may be discretized into bins of width $\delta z$.

Residue characters decompose the operator into blocks
\[
  \mathcal L_s=\mathcal L_{s,0}\oplus\bigoplus_{\chi\ne1}\mathcal L_{s,\chi}.
\]

\subsection{Test IV.1: Spectral gap as a function of tilt $s$}

\paragraph{Goal.} Determine whether H23a-type character gaps persist through the critical tilt $s=1$.

\paragraph{Procedure.}
For each $Q$, $L$, and $s\in[0,1]$:
\begin{enumerate}
  \item Build the finite matrix approximation to $\mathcal L_s$.
  \item Compute principal spectral radius $\rho_0(s)$.
  \item For each nontrivial character block, compute $\rho_\chi(s)$.
  \item Record gap
  \[
    \kappa_Q(s)=1-\max_{\chi\ne1}\frac{\rho_\chi(s)}{\rho_0(s)}.
  \]
\end{enumerate}

\paragraph{Output CSV.}
\begin{verbatim}
Q,L,z_bins,s,rho0,max_rho_chi,kappa,argmax_chi,iterations,residual
\end{verbatim}

\paragraph{Expected signal.}
A positive gap over $s\in[0,1]$ supports stability of the certificate under critical tilting.

\subsection{Test IV.2: Doob-transformed operator paths}

\paragraph{Goal.} Compare actual critical excursions to paths typical under the Doob-transformed critical operator.

\paragraph{Procedure.}
Compute principal eigenfunction $h_s$ of $\mathcal L_s$ and define
\[
  \mathcal L_s^h f=\frac{1}{\rho_s h_s}\mathcal L_s(h_sf).
\]
Sample or enumerate typical paths under $\mathcal L_1^h$ and compare valuation histograms and boundary layers to actual critical excursions.

\subsection{Test IV.3: Post-exit operator spectrum}

\paragraph{Goal.} Test the spectral form of post-exit taboo.

\paragraph{Object.} Let $\mathcal E$ be an empirical or exact post-exit kernel from completed critical excursions to next launch floors/residue-height states.

Compute spectral radius of
\[
  \mathcal K_A\mathcal E\mathcal K_A
\]
where $\mathcal K_A$ is the critical-entrance or critical-tilted operator. In the strongest case, the matrix is zero on the exact critical sector; otherwise one wants spectral radius $<1$.

\paragraph{Output CSV.}
\begin{verbatim}
Q,L,A,operator_type,rho_K,rho_KEK,max_entry,nonzero_edges,
spectral_gap,iterations,residual
\end{verbatim}

\paragraph{Expected signal.}
For $A=6$, existing tests suggest no immediate second critical event. The spectral test should show either zero transition or a strong gap.

\subsection{Test IV.4: Dirac-to-QSD comparison}

\paragraph{Goal.} Quantify the annealed-to-quenched gap.

\paragraph{Procedure.}
Start from Dirac states corresponding to actual launch floors. Iterate the annealed killed operator and compare:
\begin{enumerate}
  \item annealed distribution at time $n$;
  \item QSD profile;
  \item actual deterministic path state at time $n$.
\end{enumerate}
Record discrepancies in total variation, residue Fourier modes, and height coordinate.

\section{Auxiliary Tests for $C_A$ and Height Law}

\subsection{Critical-floor census}

Compute $C_A$ for octaves and thresholds. Output:
\begin{verbatim}
octave,A,total_odd,count_critical,count_per_octave,
min_V,max_V,min_amin,max_height,mean_word_len,mean_runs
\end{verbatim}

\subsection{Valuation-word grammar diagnostics}

For each critical word, record:
\begin{itemize}
  \item total word length;
  \item number of $b=1$ runs;
  \item longest $b=1$ run;
  \item drop-size counts for $b=2,3,4,\ldots$;
  \item entry/bulk/exit layer histograms.
\end{itemize}
This verifies that the target is not a finite grammar.

\subsection{Pressure fit from observed words}

For each octave and criticality threshold, estimate empirical transform
\[
  \widehat M_N(s)=\sum_b\widehat p_{b,N}2^{s(\log_2 3-b)}.
\]
Compare with
\[
  M(s)=\frac{3^s}{2^{1+s}-1}.
\]

\section{Exactness, Validation, and Failure Modes}

\subsection{Exact predicates versus diagnostics}

The following must be exact:
\begin{itemize}
  \item accelerated step for all tested $x$;
  \item $b=\ord(3x+1)$;
  \item peak $P$ within the integer range;
  \item criticality $V^2\le2^AP$;
  \item first return below launch floor;
  \item cycle/terminal detection if used.
\end{itemize}

The following may be floating-point diagnostics:
\begin{itemize}
  \item $h=\log_2(P/V)$;
  \item $\mathrm{amin}=\log_2(V^2/P)$;
  \item entropy and KL divergence;
  \item spectral radii and eigenvectors;
  \item pressure fits.
\end{itemize}

\subsection{Validation tests}

Every executable should include a small deterministic test suite:
\begin{enumerate}
  \item Check known short orbits entering $1$.
  \item Check $v_2(3x+1)$ for hand-computable values.
  \item Check exact criticality predicate against floating diagnostic for small values.
  \item Check histograms sum to word length.
  \item Check $M(1)=1$ numerically.
  \item Check $\sum_{b=1}^{B}3\cdot4^{-b}+4^{-B}=1$ for overflow tail.
\end{enumerate}

\subsection{Interpretation failure modes}

The computations may falsify specific hypotheses:
\begin{itemize}
  \item If critical bulk words do not approach $p^{(1)}$, the Cram\'er-tilt model is incomplete.
  \item If post-exit windows also approach $p^{(1)}$, the taboo mechanism is not explained by anti-tilt.
  \item If $\rho(\mathcal K_A\mathcal E\mathcal K_A)\ge1$, the spectral taboo form fails.
  \item If character gaps collapse near $s=1$, H23a is not stable under critical tilt.
  \item If Dirac-to-QSD discrepancy remains uncontrolled, the annealed-to-quenched gap remains exactly where expected.
\end{itemize}

\section{Output File Formats}

\subsection{Raw excursion records}

\begin{verbatim}
V,octave,A,critical,steps,peak,exit_floor,height,amin,
word_len,max_b,b_word_hash,b1_runs,max_b1_run,overflow_flag
\end{verbatim}

The full word may be optionally stored in a separate compressed text/binary file keyed by \texttt{b\_word\_hash}.

\subsection{Histogram summaries}

\begin{verbatim}
group_id,label,octave,A,region,count_words,count_symbols,b,count_b,
p_hat,p_annealed,p_tilted
\end{verbatim}

\subsection{Entropy-space windows}

\begin{verbatim}
x0,n,window,label,mean_b,mean_delta,entropy,
kl_to_p,kl_to_tilt,best_s,kl_to_exp_family,height_start,height_end
\end{verbatim}

\subsection{Spectral summaries}

\begin{verbatim}
Q,L,z_bins,s,block,rho,residual,iterations,
rho0,max_rho_chi,kappa,argmax_chi
\end{verbatim}

\section{Milestones}

\subsection*{Milestone 1: Exact core and critical census}
\begin{itemize}
  \item Implement accelerated step, excursion extraction, exact criticality.
  \item Reproduce known $C_6$ counts over a small range.
  \item Produce octave count table for $A=4,5,6,7,8$.
\end{itemize}

\subsection*{Milestone 2: Doob/Cram\'er diagnostics}
\begin{itemize}
  \item Compute critical bulk valuation histograms.
  \item Verify convergence toward $p_b^{(1)}=3\cdot4^{-b}$.
  \item Fit height-tail law $\Prob(h\ge H)\asymp2^{-H}$.
\end{itemize}

\subsection*{Milestone 3: Bridge and post-exit diagnostics}
\begin{itemize}
  \item Split critical words into entry/bulk/exit layers.
  \item Measure post-exit law and $M_{post}(1)$.
  \item Test $E(C_6)\cap C_6=\varnothing$ at increased scale.
\end{itemize}

\subsection*{Milestone 4: Entropy-space transport}
\begin{itemize}
  \item Generate sliding-window entropy trajectories.
  \item Fit best exponential-family parameter $s$.
  \item Test whether post-exit windows avoid neighborhoods of $s=1$.
\end{itemize}

\subsection*{Milestone 5: Transfer-operator flow}
\begin{itemize}
  \item Build finite residue-height operators $\mathcal L_s$.
  \item Compute character gaps over $s\in[0,1]$.
  \item Construct empirical/exact post-exit kernel $\mathcal E$ and compute $\rho(\mathcal K_A\mathcal E\mathcal K_A)$.
\end{itemize}

\section{Deliverables}

Each round should produce:
\begin{enumerate}
  \item source code and build log;
  \item raw CSV files;
  \item summary CSV/JSON files;
  \item a short report with plots and interpretation;
  \item exact command lines sufficient to reproduce the run;
  \item a manifest of tested ranges and overflow exclusions.
\end{enumerate}

\section{Conclusion}

The computation program should test a single organizing hypothesis from four directions:
\begin{quote}
A deterministic branch that shadows the annealed rare-event flow for too long must force a finite spectral/character obstruction, or else reveal a new post-exit taboo mechanism in the height cocycle.
\end{quote}
The Doob/Cram\'er tests identify the local rare-event law. The Schr\"odinger-bridge tests isolate entry and exit boundary layers. The entropy-transport tests track deterministic empirical laws in information-geometric coordinates. The transfer-operator tests connect the same phenomenon to H23a, QSD, and post-exit spectra. Together they are designed to distinguish three possibilities: a provable taboo transition, a finite spectral obstruction, or a genuinely new quenched equidistribution barrier.

\end{document}
